\section{Preliminaries}\label{sec:prelim}

All graphs are finite and simple. For a graph $G$ and a vertex $v$ we write $\Nb(v)=\Nb_G(v)$ for the \emph{open} neighbourhood (the subgraph induced by the vertices adjacent to $v$), and for distinct vertices $u,v$ we write $\Nb(u,v)=\Nb(u)\cap\Nb(v)$ for the \emph{common neighbourhood}, again as an induced subgraph. $G-v$ and $G-uv$ denote deletion of a vertex and of an edge; $G+uv$ denotes the addition of an edge between non-adjacent vertices. The \emph{clique complex} $\cl(G)$ is the simplicial complex whose simplices are the cliques of $G$; a complex of this form is called \emph{flag}. Note that $\cl(\Nb(v))$ is the link of $v$ in $\cl(G)$ and $\cl(\Nb(u,v))$ is the link of $uv$ when $uv$ is an edge.

\subsection{Contractible transformations}

We follow the conventions of~\cite{DEFH2023}.

\begin{definition}[\cite{Ivashchenko1994a}]\label{def:I}
The class $\I$ of \emph{I-contractible} graphs is the smallest class of graphs containing $K_1$ and closed under the following four transformations:
\begin{enumerate}[label=(I\arabic*),leftmargin=3em]
\item deleting a vertex $v$ with $\Nb(v)\in\I$;
\item adding a new vertex $v$ adjacent to the vertices of an induced subgraph $H\in\I$;
\item deleting an edge $uv$ with $\Nb(u,v)\in\I$;
\item adding an edge $uv$ between non-adjacent vertices with $\Nb(u,v)\in\I$.
\end{enumerate}
\end{definition}

The four transformations are the \emph{contractible transformations} of Ivashchenko; (I3) and (I4) can be realised by (I1) and (I2)~\cite[Lemma~3]{DEFH2023}. Chen, Yau and Yeh~\cite{ChenYauYeh2001} proved that these transformations preserve the simple homotopy type of the clique complex. Hence every graph in $\I$ has contractible $\cl(G)$; the converse is~\cite[Corollary~1.5]{DochtermannMatsushita2025}.

\begin{definition}[\cite{DEFH2023}]\label{def:IS}
The class $\IS$ of \emph{strong I-contractible} graphs is defined recursively as follows. The graph $K_1$ belongs to $\IS$; a graph $G$ with at least two vertices belongs to $\IS$ if at least one of the following conditions holds:
\begin{enumerate}[label=(S\arabic*),leftmargin=3em]
\item $G-v\in\IS$ and $\Nb_G(v)\in\IS$ for some vertex $v$;
\item $G-uv\in\IS$ and $\Nb_G(u,v)\in\IS$ for some edge $uv$.
\end{enumerate}
The empty graph does not belong to $\IS$. Equivalently, $G\in\IS$ if it has a sequence of these deletions ending at $K_1$. The class $\IVS$ of \emph{vertex-strong I-contractible} graphs is defined by the same recursion using only vertex deletions and requiring both $G-v$ and $\Nb_G(v)$ to belong to $\IVS$.
\end{definition}

The empty graph is \emph{not} in $\IS$; in particular an edge $uv$ with $\Nb(u,v)=\emptyset$ is never deleted, and an isolated vertex is never deleted. A vertex $v$ with $\Nb(v)\in\IS$, or an edge $uv$ with $\Nb(u,v)\in\IS$, will be called \emph{$\IS$-deletable}. In a graph of $\IS$ the order of deletions matters: a sequence may lead to a graph with no $\IS$-deletable element while another sequence reaches $K_1$ (\cite[Fig.~2B]{DEFH2023}); thus a greedy procedure does not decide $\IS$ and an exact search is needed (Section~\ref{sec:algorithms}).

\begin{definition}[\cite{DEFH2023,FieuxJouve2020}]\label{def:kdis}
A vertex $v$ of $G$ is \emph{$0$-dismantlable} if $\Nb(v)$ is a cone (some vertex of $\Nb(v)$ is adjacent to all others), and, for $k\ge 1$, \emph{$k$-dismantlable} if $\Nb(v)$ is $(k-1)$-dismantlable; $(-1)$-dismantlable graphs are the cones. A graph is \emph{$k$-dismantlable} ($k$-Dis) if it can be reduced to $K_1$ by deleting $k$-dismantlable vertices. A $0$-dismantlable graph is also called \emph{dismantlable} or \emph{cop-win}~\cite{NowakowskiWinkler1983}.
\end{definition}

One has $(-1)\text{-Dis}\subseteq 0\text{-Dis}\subseteq 1\text{-Dis}\subseteq\cdots\subseteq\IVS\subseteq\IS\subseteq\I$ (\cite[Lemma~19]{DEFH2023} and induction).

\subsection{Homology}

We write $\widetilde H_i(G;R)$ for the reduced simplicial homology of $\cl(G)$ with coefficients in a ring $R$, and $\chi(G)$ for the Euler characteristic of $\cl(G)$, which can be computed from the clique numbers $c_k(G)$ as $\chi(G)=\sum_{k\ge1}(-1)^{k-1}c_k(G)$. A graph is \emph{$\Ztwo$-acyclic} if it is non-empty and $\widetilde H_i(G;\Ztwo)=0$ for every $i\ge0$; in particular it is connected. Define \emph{$\ZZ$-acyclicity} similarly. The universal coefficient theorem implies that $\ZZ$-acyclicity implies $\Ztwo$-acyclicity. For a finite connected complex, $\Ztwo$-acyclicity is equivalent to all its positive-degree integral homology groups being finite of odd order; integral acyclicity additionally requires these groups to vanish. A simplicial complex is \emph{contractible} if it is homotopy equivalent to a point. Contractibility implies integral acyclicity, but the converse is false in general: an acyclic complex can have a non-trivial perfect fundamental group. We establish contractibility of our examples by explicit I-transformations, rather than inferring it from integral homology alone.


Integer homology of the small complexes in this paper is computed by exact Smith normal form of the boundary matrices (Section~\ref{sec:algorithms}). We use the following elementary fact repeatedly.

\begin{lemma}\label{lem:acyclic-IS}
Let $n$ be the smallest order of a $\Ztwo$-acyclic graph that is not in $\IS$ (if such graphs exist). Then for every graph $H$ with fewer than $n$ vertices, $H\in\IS$ if and only if $H$ is $\Ztwo$-acyclic.
\end{lemma}
\begin{proof}
Graphs in $\IS$ are in $\I$, hence have contractible clique complex and are $\Ztwo$-acyclic; the converse holds below $n$ by minimality of $n$.
\end{proof}

\subsection{Discs with identified boundary}

For the structural descriptions in Section~\ref{sec:examples}, let $X$ be a pure $2$-dimensional simplicial complex. Start with disjoint copies of its triangles and glue corresponding edges only when the original edge belongs to exactly two triangles. Identify the corresponding corners under these gluings. This produces a complex $\widetilde X$ with a natural map to $X$; in each case discussed below, we check that $\widetilde X$ is a compact surface with boundary. Thus the construction cuts along every edge whose triangle incidence differs from two, including edges in four triangles. Purity alone does not guarantee that the result is a surface. For the stuck $K_4$-free examples every edge belongs to at least two triangles, so the cut occurs precisely along the singular edges.

Suppose that $\widetilde X$ is a disc, the map is injective on its interior, and $X$ is recovered solely by identifying boundary points according to their image. Let $S$ be the image of the boundary, and suppose that $S$ is a connected graph with exactly one cycle $C$. Then $X$ is the adjunction space $S\cup_f D^2$, where the boundary word records the attaching map $f$. Contract the trees attached to $C$ and orient $C$. If the resulting boundary map has degree $d$, then
\[
 X\simeq S^1\cup_d D^2.
\]
Indeed, collapsing a maximal tree in $S$ is a homotopy equivalence of the adjunction spaces, and maps $S^1\to S^1$ are classified up to homotopy by their degree. In particular, $|d|=1$ gives a contractible complex, while $|d|=3$ gives the Moore space $M(\ZZ/3,1)$ up to homotopy. When the boundary makes three complete passes around $C$, the patterns $a\,a\,a^{-1}$ and $a^3$ have degrees $1$ and $3$, respectively. Excursions into an attached tree contribute degree zero.

These hypotheses are verified for $\Gt$, the $15$-vertex dunce hat $\Dunce$ of~\cite{DEFH2023}, and $48$ of the $51$ stuck $K_4$-free graphs on $14$ vertices. Among the latter, three have degree $\pm1$ and $45$ have degree $\pm3$. The remaining three unfold into two annuli and one disconnected surface; the single-disc argument does not apply to them.
